\documentclass[UTF8]{ctexart}
\usepackage{amsmath}
\usepackage{booktabs}
\usepackage{geometry}
\usepackage{listings}
\usepackage{xcolor}
\usepackage{hyperref}
\usepackage{graphicx}

\geometry{a4paper, left=2.5cm, right=2.5cm, top=2.5cm, bottom=2.5cm}
\hypersetup{colorlinks=true, linkcolor=blue, urlcolor=blue, citecolor=blue}

\lstset{
  basicstyle=\ttfamily\small,
  breaklines=true,
  frame=single,
  keywordstyle=\color{blue},
  commentstyle=\color{gray},
  stringstyle=\color{teal},
  showstringspaces=false,
  tabsize=2,
  columns=flexible,
}

\title{\textbf{系统工具开发课程实验报告（第四周）}}
\author{廖世嘉 \quad 25020007073 \\ 25级计算机科学与技术}
\date{2026年9月14日}

\begin{document}
\maketitle
\tableofcontents
\newpage

\section{实验基本信息}

\begin{table}[htbp]
\centering
\caption{实验基本信息}
\label{tab:info}
\begin{tabular}{ll}
\toprule
项目 & 内容 \\
\midrule
姓名 & 廖世嘉 \\
学号 & 25020007073 \\
专业 & 25级计算机科学与技术 \\
实验环境 & Linux / Python / ruff / pytest / Make / curl / jq / Git \\
实验日期 & 2026年9月14日 \\
\bottomrule
\end{tabular}
\end{table}

\section{实验目的}

本实验围绕代码质量门禁、Make 构建系统、API 数据处理与报告生成、以及综合仓库交付四个主题展开。通过四个连续实验，建立本地质量检查流程，掌握 Make 的增量构建机制，使用 curl 与 jq 处理 API 数据并生成 Markdown 报告，以及完成一次包含测试、构建、哈希计算和 Git 提交的综合检查。

\section{实验环境与工具}

本实验在 Linux 和 Windows 环境中完成，主要使用 ruff（代码格式化与检查）、pytest（测试框架）、Make（构建系统）、curl（HTTP 请求）、jq（JSON 处理）、python -m http.server（本地 HTTP 服务）、python -m build（wheel 构建）、sha256sum（哈希计算）以及 Git（版本控制）。实验目录为 SystemToolkitDevCourse，四个实验分别位于 q13、q14、q15 和 q16。

\section{实验十三：建立可执行的本地质量门禁}

\subsection{实验要求}

复制 q10 为 q13，在 pyproject.toml 中加入 ruff 和 pytest 的最小配置。补充一个正常姓名测试和一个空白姓名测试，每个测试包含有效断言。运行 ruff format、ruff check 和 pytest，修复全部问题，不得全局忽略规则。编写 check.sh，依次执行 ruff format --check、ruff check 和 pytest，运行并保证返回 0。

\subsection{配置文件}

pyproject.toml 中新增 ruff 和 pytest 配置：

\begin{lstlisting}
[tool.ruff]
line-length = 88
target-version = "py310"

[tool.ruff.lint]
select = ["E", "F", "W", "I"]

[tool.pytest.ini_options]
pythonpath = ["src"]
testpaths = ["."]
\end{lstlisting}

\subsection{测试用例}

test\_cli.py 包含两个测试：

\begin{lstlisting}[language=Python]
import pytest
from greetlab.cli import main

def test_normal_name(capsys):
    import sys
    sys.argv = ["sdt-greet", "--name", "alice"]
    main()
    out = capsys.readouterr().out
    assert "Hello, alice!" in out

def test_blank_name_exits_nonzero():
    import sys
    sys.argv = ["sdt-greet", "--name", "   "]
    with pytest.raises(SystemExit) as exc_info:
        main()
    assert exc_info.value.code == 2
\end{lstlisting}

\subsection{质量检查脚本}

check.sh 依次执行三项检查：

\begin{lstlisting}[language=bash]
#!/usr/bin/env bash
set -euo pipefail

echo "=== ruff format --check ==="
ruff format --check .

echo "=== ruff check ==="
ruff check .

echo "=== pytest ==="
python -m pytest test_cli.py -v

echo "=== All checks passed ==="
\end{lstlisting}

\subsection{结果}

运行 check.sh 输出：

\begin{lstlisting}
=== ruff format --check ===
4 files already formatted
=== ruff check ===
All checks passed!
=== pytest ===
2 passed
=== All checks passed ===
\end{lstlisting}

退出码为 0，全部检查通过。期间 ruff 发现的导入排序问题已通过 ruff check --fix 修复，未使用全局忽略规则。

\section{实验十四：让 Make 只重建真正受影响的产物}

\subsection{实验要求}

在 q14 中创建 data.csv、stats.py、build\_report.py 和 report.md 四个源文件，编写 Makefile 包含 all、stats.txt、report.txt 和 clean 目标。准确列出数据文件、脚本和中间产物依赖，将 all 和 clean 声明为 .PHONY。验证首次 make 生成两个产物，第二次无改动时不执行配方，touch data.csv 后重新生成，clean 只删除生成物。

\subsection{源文件}

\begin{lstlisting}[language=bash]
# data.csv
name,value
alpha,2
beta,3
gamma,5
\end{lstlisting}

stats.py 读取 data.csv 并计算总和写入 stats.txt：

\begin{lstlisting}[language=Python]
import csv
with open("data.csv") as f:
    total = sum(int(r["value"]) for r in csv.DictReader(f))
open("stats.txt", "w").write(str(total))
\end{lstlisting}

build\_report.py 将 report.md 和 stats.txt 合并为 report.txt：

\begin{lstlisting}[language=Python]
text = open("report.md").read()
total = open("stats.txt").read()
open("report.txt", "w").write(f"{text}\nTotal: {total}\n")
\end{lstlisting}

\subsection{Makefile}

\begin{lstlisting}
.PHONY: all clean

all: report.txt

stats.txt: data.csv stats.py
	python3 stats.py

report.txt: report.md stats.txt build_report.py
	python3 build_report.py

clean:
	rm -f stats.txt report.txt
\end{lstlisting}

依赖链：data.csv 和 stats.py 生成 stats.txt，report.md、stats.txt 和 build\_report.py 生成 report.txt。all 依赖 report.txt，clean 删除两个生成物。

\subsection{验证结果}

\begin{lstlisting}[language=bash]
# 首次 make
$ make
python3 stats.py
python3 build_report.py

# 第二次 make（无改动）
$ make
make: Nothing to done for 'all'.

# touch data.csv 后
$ touch data.csv && make
python3 stats.py
python3 build_report.py

# clean
$ make clean
rm -f stats.txt report.txt
\end{lstlisting}

首次 make 生成两个产物，第二次 make 因无改动而不执行配方，touch data.csv 后 stats.txt 和 report.txt 均重新生成，clean 只删除生成物，保留源文件。Make 的增量构建机制工作正确。

\section{实验十五：把本地 API 数据转换为可读报告}

\subsection{实验要求}

在 q15 中创建给定的 packages.json，运行 python -m http.server 8000 启动本地 HTTP 服务。使用 curl -fsS 获取数据，使用 jq 筛选 status 为 active 且 downloads 不少于 100 的记录，按 downloads 降序、name 升序排列。编写 api\_report.sh 将结果生成 summary.md，包含标题和 name/version/downloads 三列表格。

\subsection{数据文件}

packages.json 包含 5 条记录：

\begin{lstlisting}[language=json]
[
  {"name":"alpha","status":"active","downloads":120,"version":"1.2.0"},
  {"name":"beta","status":"inactive","downloads":900,"version":"2.0.0"},
  {"name":"gamma","status":"active","downloads":450,"version":"1.5.1"},
  {"name":"delta","status":"active","downloads":450,"version":"0.9.0"},
  {"name":"epsilon","status":"active","downloads":80,"version":"3.1.0"}
]
\end{lstlisting}

\subsection{脚本实现}

api\_report.sh 使用 curl 获取 JSON 数据，通过 jq 进行筛选、排序和 Markdown 表格生成：

\begin{lstlisting}[language=bash]
#!/usr/bin/env bash
set -euo pipefail

URL="http://127.0.0.1:8000/packages.json"
OUTPUT="summary.md"

curl -fsS "$URL" | \
  jq -r '
    ["| name | version | downloads |", "|------|---------|----------|"] as $header |
    [$header[],
     ([.[] | select(.status == "active" and .downloads >= 100)]
      | sort_by(-.downloads, .name)
      | .[]
      | "| \(.name) | \(.version) | \(.downloads) |")
    ] | .[]' > "$OUTPUT"

{
  echo "# Package Summary"
  echo ""
  cat "$OUTPUT"
} > "${OUTPUT}.tmp" && mv "${OUTPUT}.tmp" "$OUTPUT"

echo "Report written to $OUTPUT"
\end{lstlisting}

\subsection{结果}

在 q15 目录启动 HTTP 服务后运行 api\_report.sh，生成的 summary.md：

\begin{lstlisting}
# Package Summary

| name | version | downloads |
|------|---------|----------|
| delta | 0.9.0 | 450 |
| gamma | 1.5.1 | 450 |
| alpha | 1.2.0 | 120 |
\end{lstlisting}

筛选结果正确：beta（inactive）和 epsilon（downloads=80 < 100）被排除。delta 和 gamma 的 downloads 均为 450，按 name 升序 delta 排在 gamma 之前。curl -fsS 在服务不可用时返回非零退出码，确保错误可被脚本捕获。

\section{实验十六：修复并交付一个陌生的小型工具仓库}

\subsection{实验要求}

复制 q13 为 q16 并初始化 Git。把问候语实现临时改成始终返回字面量 Hello, name!，运行 check.sh 确认测试失败，随后定位并修复。编写最小 Makefile 包含 check、build 和 clean，check 调用 check.sh，build 生成 wheel。运行 make check 和 make build，计算 wheel 的 SHA-256，并把最终改动提交为一个内容聚焦的 Git 提交。

\subsection{临时破坏与验证}

将 cli.py 中的实现临时改为字面量输出：

\begin{lstlisting}[language=Python]
# 临时改坏
print("Hello, name!")
\end{lstlisting}

运行 check.sh 确认测试失败：

\begin{lstlisting}
=== ruff check ===
F841 Local variable `a` is assigned to but never used
=== pytest ===
FAILED test_normal_name - assert "Hello, alice!" in "Hello, name!"
FAILED test_blank_name_exits_nonzero - Failed: DID NOT RAISE SystemExit
\end{lstlisting}

\subsection{修复}

将实现恢复为正确版本：

\begin{lstlisting}[language=Python]
if not a.name.strip():
    raise SystemExit(2)
print(f"Hello, {a.name}!")
\end{lstlisting}

\subsection{Makefile}

\begin{lstlisting}
.PHONY: check build clean

check:
	bash check.sh

build:
	python3 -m build --no-isolation

clean:
	rm -rf build dist src/*.egg-info .pytest_cache .ruff_cache
\end{lstlisting}

\subsection{构建与交付}

运行 make check 全部通过，运行 make build 成功生成 wheel。计算 SHA-256：

\begin{lstlisting}[language=bash]
$ sha256sum dist/*.whl
f5acb7abe8a6604d88a68cedc447857cf9b7addd89e04c4a154ca489b630d1c6  dist/greetlab_25020007073-0.1.0-py3-none-any.whl
\end{lstlisting}

Git 提交记录：

\begin{lstlisting}[language=bash]
$ git log --oneline
778d461 fix: restore correct greeting and add Makefile with check/build/clean
c4cd828 init: greetlab package from q13
\end{lstlisting}

提交信息包含问题说明、修复内容和 SHA-256 哈希值，内容聚焦于本次修复和构建配置。

\section{实验结果汇总}

\begin{table}[htbp]
\centering
\caption{四个实验的完成情况}
\label{tab:summary}
\begin{tabular}{cll}
\toprule
题号 & 实验主题 & 验证结果 \\
\midrule
13 & 代码质量门禁 & ruff format/check 和 pytest 全部通过，check.sh 返回 0 \\
14 & Make 构建系统 & 首次构建、无改动跳过、touch 重建、clean 清理均正确 \\
15 & API 数据转报告 & curl+jq 筛选排序正确，summary.md 表格生成成功 \\
16 & 综合测试 & 临时破坏验证失败，修复后 make check/build 通过，SHA-256 已计算并提交 \\
\bottomrule
\end{tabular}
\end{table}

\section{问题与解决方法}

\subsection{ruff 导入排序问题}

在实验十三中，初次运行 ruff check 发现 test\_cli.py 的导入顺序不符合 isort 规则（I001）。由于 pyproject.toml 中启用了 I 规则，ruff 要求导入按标准库、第三方库、本地模块的顺序排列。使用 ruff check --fix 自动修复后，再次检查全部通过，无需全局忽略规则。

\subsection{Makefile 中的 Tab 缩进}

在实验十四中，Makefile 的配方行必须使用 Tab 而非空格缩进。在 Windows 编辑器中保存时，有时会将 Tab 转换为空格，导致 make 报错"missing separator"。解决方法是在 WSL 中使用 cat -A 检查 Makefile 中的缩进字符，确保配方行以 \texttt{\textasciitab} 开头。

\subsection{jq 排序中的多字段排序}

在实验十五中，需要先按 downloads 降序、再按 name 升序排列。jq 的 sort\_by 函数默认按升序排列，对于降序需要在字段前加负号：sort\_by(-.downloads, .name)。这一语法确保 delta 和 gamma 在相同 downloads（450）时按 name 升序排列。

\subsection{Git 提交的聚焦性}

在实验十六中，初始提交误将临时文件 commit\_msg.txt 一并提交。通过 git rm --cached 移除该文件并 amend 提交，确保最终提交仅包含 .gitignore 和 Makefile 两个必要文件的变更，符合"内容聚焦"的要求。

\section{实验总结}

通过本次实验，我建立了完整的代码质量保障流程：从 ruff 的格式化和静态检查，到 pytest 的自动化测试，再到 check.sh 的一键执行，形成了可重复的质量门禁。Make 构建系统的实验让我深入理解了文件依赖关系和增量构建的本质——只有当依赖比目标更新时才执行配方，这避免了不必要的重复构建。

实验十五让我掌握了 curl 与 jq 的组合使用方式：curl -fsS 确保请求失败时脚本中断，jq 的管道和 select/sort\_by 函数可以高效地处理 JSON 数据。实验十六作为综合测试，将前序所有技能整合到一个可交付的 Git 仓库中，从临时破坏到修复、从构建到哈希计算、从测试到提交，形成了一次完整的"红——绿——交付"循环。

四周实验的核心收获是：好的工程实践不在于单个工具的使用，而在于将质量检查、构建系统、版本控制和自动化测试组合成一个可重复、可验证的工作流。check.sh 和 Makefile 就是这种工作流的物理载体——它们将隐性的操作步骤固化为显性的、可执行的命令序列。

\end{document}
